% 使用 XeLaTeX 编译；正文控制为一页 A4。
\documentclass[UTF8,a4paper,11pt,fontset=windows]{ctexart}
\usepackage[left=2.2cm,right=2.2cm,top=1.8cm,bottom=1.8cm]{geometry}
\usepackage{amsmath}
\usepackage{microtype}
\usepackage[hidelinks]{hyperref}
\pagestyle{empty}
\setlength{\parindent}{2em}
\setlength{\parskip}{0.45em}
\linespread{1.10}
% 小标题独占一行，正文另起一段，并避免标题单独留在页末。
\newcommand{\lead}[1]{\par\noindent\textbf{#1}\par\nobreak}

\begin{document}

\begin{center}
  {\zihao{3}\bfseries AI辅助高中数学实践建模项目计划}\par
  \vspace{0.3em}
  {\zihao{5} 林立康，计算数学，25210180078}
\end{center}
\vspace{-0.5em}

\lead{一、项目构想与研究问题}
本项目回应高中数学课标关于开展数学建模活动、经历完整建模过程的要求，拟开发一个AI辅助的实践项目设计网页。设计包括两个相互配合的部分：一是根据项目需要与学生基础，重组相关数学知识的学习顺序；二是围绕一个较大型的真实问题，将完整项目拆分为前后连贯、逐步深入的小问和实践任务。学生通过搜集数据、提出假设、选择模型、计算检验和修改方案，逐步完成项目。每一阶段形成的数据、图表或模型都用于后续任务，并为后续知识学习作准备。核心研究问题是：\textbf{AI如何将真实建模项目转化为高中生可以逐步开展的实践任务，并合理安排所需知识？任务分解能否兼顾过程引导与学生自主决策，使学生有机会选择方法、解释结果和修改模型？}

\lead{二、目标使用人群与实际参与者}
\textbf{目标使用人群为参加数学建模课程、社团或研究性学习的高中学生。受限于招募中学生的实际困难，本次拟招募6至10名数学系研究生参与原型试用与内容评议。}参与者完整审阅项目方案，并试做数据处理、模型拟合与检验等关键环节，每人约60至90分钟，随后进行访谈。记录其高中数学教学经历、建模经验和工具使用经验，用于理解反馈差异。本次研究评价项目设计与实施流程的可行性，研究生反馈用于修订原型；高中生的实际学习效果留待后续课堂研究检验。

\lead{三、实践项目示例与AI支持方式}
以“利用公开天文数据拟合火星轨道”为示例，AI将项目组织为六个相互衔接的阶段，每阶段再细分为若干引导性小问：第一，明确要描述的轨道及研究时段，查找并下载位置数据，记录来源、单位与生成方式；第二，统一日心坐标与单位，在近似轨道平面上绘制散点图，说明简化假设；第三，根据数据形态提出候选模型，比较圆与椭圆模型的选择依据；第四，借助电子表格或编程工具估计参数，解释参数的几何意义；第五，用预留的位置数据检验拟合误差，分析误差来源并修订模型；第六，整理数据、图表和模型文件，撰写报告，说明方法选择、结果及适用范围。拟采用JPL Horizons公开历表，注明其为计算生成的位置数据。AI根据各阶段需要安排向量、圆锥曲线和数据分析等知识的学习顺序，标明前一阶段成果如何用于下一阶段；同时提供操作提示和思路反馈，逐步减少帮助，将数据筛选、模型选择和改进理由留给学生判断。

\lead{四、关键条件与评价设计}
重点考虑\textbf{学生已有数学知识、数据与工具的可获得性、项目所需时间、提示详细程度和教师支持}。试用者对照同一项目的整体任务说明与AI生成的分阶段方案，检查知识安排是否合理、操作是否可完成、阶段成果是否衔接，以及是否保留提出假设、比较模型和修订方案的空间。评价证据包括数据来源记录、拟合图表、模型修改说明、网页操作记录和访谈反馈；重点分析任务中断的原因、提示能否帮助继续实践，以及哪些步骤引导过细或缺少支持。研究以问题清单、评分分布和修改前后案例呈现结果，为高中课堂中的进一步试用形成具体建议。

\lead{五、实施要求与预期成果}
研究者核验数据来源、数学内容与工具操作，保留AI初稿和修订记录；试用取得参与者知情同意，匿名记录反馈。预期成果包括一个网页原型、一套包含真实数据来源、递进实践任务、知识学习顺序和评价标准的高中数学建模项目，以及一份原型试用报告。

\vspace{0.2em}
{\footnotesize\noindent 设计依据：\href{https://jyj.changdu.gov.cn/cdsjyj/c101476/202205/764ee1914e3c45aaba8e372f1b628a34/files/75676dfa9af0442db1457404a884a69f.pdf}{《普通高中数学课程标准（2017年版2020年修订）》第34至36页}。数据来源说明：\href{https://ssd.jpl.nasa.gov/horizons/manual.html}{JPL Horizons官方文档}。\par}

\end{document}
